% -------------------------------------------------
% VERSION TRACKER — No modificar este bloque
% Versión:    Tesis Humanizada
% Archivo:    Capitulo1Introduccion/Objetivos.tex
% Última mod:  2026-05-08T08:16:48
% Git commit:  cc9c766f @ main
% Audiencia:   general
% Verificar antes de editar: bash check_tesis_version.sh overleaf_tesis_humanizado/Capitulo1Introduccion/Objetivos.tex
% -------------------------------------------------


\section{Objetivos}

\subsection{Objetivo general}

Implementar un sistema web con arquitectura moderna y un módulo de consulta inteligente, que ayude a las PYMEs colombianas a optimizar el control de su bodega, reducir pérdidas por vencimiento y tomar mejores decisiones —dejando atrás los procesos manuales, la falta de trazabilidad y la ausencia de reportes.

\subsection{Objetivos específicos}

Cada objetivo incluye su indicador de cumplimiento para que el lector pueda verificar si se alcanzó lo propuesto.

\begin{enumerate}

\item \textbf{Entender las necesidades reales de una PYME.} Analizar cómo funciona la bodega de una panificadora real del Valle del Cauca, mediante entrevistas guiadas con el personal, para identificar y documentar al menos 9 casos de uso, 10 reglas de negocio y los criterios que debe cumplir el sistema.

\textit{Indicador:} Documento de requisitos (estándar IEEE 830) con 9 casos de uso validados por la empresa caso de estudio y aprobados por el director del trabajo de grado.

\item \textbf{Diseñar el plano del sistema.} Crear el modelo de la base de datos (mínimo 10 tablas), los diagramas UML de los 4 procesos más importantes (entrada, despacho, conciliación y consulta inteligente) y al menos 18 prototipos visuales (mockups) en Figma que muestren cómo se verá cada pantalla. Incluir el diseño del módulo de consulta inteligente con su clasificador de preguntas en 10 categorías.

\textit{Indicador:} Documento de diseño con modelo entidad-relación de 10 tablas, diagramas UML para 4 procesos críticos, 18 mockups en Figma aprobados por el director y especificación del motor RAG con 10 categorías de clasificación de intención.

\item \textbf{Construir el sistema por etapas.} Desarrollar los 9 módulos del sistema usando la metodología ágil SCRUM en ciclos de dos semanas. Implementar la gestión de materiales y lotes, el despacho automático por fecha de vencimiento (FEFO), el historial inmutable de movimientos (Kardex), los roles de usuario (Admin y Operario), la conciliación de inventario, los reportes en PDF y CSV, las etiquetas QR para identificar lotes y el asistente de consultas en lenguaje natural conectado a la base de datos real.

\textit{Indicador:} Sistema web funcional con 9 módulos operativos, repositorio en GitHub con al menos 50 commits y el sistema desplegado en al menos un servidor en la nube.

\item \textbf{Verificar que todo funcione correctamente.} Ejecutar un plan de pruebas con 13 conjuntos de casos que cubran: pruebas automáticas de los 9 casos de uso (con la herramienta Playwright), 7 escenarios de despacho FEFO, 50 consultas en español para medir la precisión del asistente inteligente, evaluación de usabilidad con los 10 principios de Jakob Nielsen y una auditoría de accesibilidad según el nivel AA de las pautas WCAG 2.1. El sistema debe alcanzar al menos un 95\% de pruebas exitosas.

\textit{Indicador:} Documento de resultados con 111 casos ejecutados, tasa de éxito global $\geq$ 95\%, informe de evaluación heurística y auditoría de accesibilidad WCAG AA.

\end{enumerate}
